
\SGAStatement{Corollary 5.2.4 bis}{With the notation of \SGARefLink{sga3:xvii:lemma:5:2:4}{5.2.4},
suppose in addition that \(U\) is flat and of finite presentation
over \(S\), with unipotent fibres, and that \(n\) is prime to the
residue characteristics of \(S\).  Then, in cases (a) and (b)
above,
\[
  H^i(H,U)=0\qquad(i>0).
\]}
\SGARefTarget{sga3:xvii:corollary:5:2:4:bis}

\SGAProof{Proof}
It is enough to show that the \(n\)-th power map on \(U\) is an
isomorphism, for it will then follow that multiplication by \(n\)
on \(H^i(H,U)\) is both an isomorphism and the zero morphism, and
hence that \(H^i(H,U)=0\).  Under the hypotheses made on \(U\), it
is enough to check that the \(n\)-th power map is an isomorphism
on the fibres of \(U\) (EGA IV, 17.9.5).  We are thus reduced to
the case in which \(S\) is the spectrum of a field \(k\) of
characteristic \(p\).  Since \((n,p)=1\), the \(n\)-th power map
on \(U\) is étale (Exposé VII); it is also a monomorphism (2.4(i)),
and hence an open immersion (EGA IV, 17.9.1).  This already shows
that its restriction to the identity component \(U^0\) is an
isomorphism.  Since \(U/U^0\) is a finite \(p\)-group, the proof
is complete.

This completes the proof of \SGARefLink{sga3:xvii:lemma:5:2:3:item:i}{5.2.3(i)}\SGARefLink{sga3:xvii:appendix:A}{(a)}, since an étale group of
multiplicative type \(H\) has order prime to \(p\).

\SGAProof{Proof of \SGARefLink{sga3:xvii:lemma:5:2:3:item:ii}{5.2.3(ii)}}
It is clear that \(R\) is a sheaf for the fpqc topology.
Moreover, in view of descent for affine schemes, the assertions
of \SGARefLink{sga3:xvii:lemma:5:2:3:item:ii}{5.2.3(ii)} are local for the fpqc topology.  We may therefore
suppose that \(k\) is algebraically closed.  The group \(H\) is
then split and the extension \(E\) is trivial by (i)\SGARefLink{sga3:xvii:section:b}{(b)}; let
\(H'\) be a lift of \(H\) in \(E\).  For every \(k\)-prescheme
\(S\) and every lift \(H''\) of \(H_S\) in \(E_S\), the groups
\(H'_S\) and \(H''\) are conjugate by an element of \(U(S)\),
since
\[
  H^1(H_S,U_S)=0
\]
by \SGARefLink{sga3:xvii:corollary:5:2:4:bis}{5.2.4 bis}.

Let \(U^{H'}\) be the sheaf of invariants of \(U\) under \(H'\).
It is representable by an algebraic subgroup of \(U\)
(\SGARefLink{sga3:viii:examples:6:5:item:d:01}{Exposé VIII, 6.5(d)}).  The preceding remarks show that the
\(k\)-morphism
\[
  U\longrightarrow R,
  \qquad
  u\longmapsto \operatorname{int}(u)H'
  \quad (u\in U(S))
\]
induces a \(k\)-isomorphism
\[
  U/U^{H'}\xrightarrow{\ \sim\ }R.
\]
This proves that \(R\) is representable and affine (\SGARefLink{sga3:xvii:proposition:2:1}{2.1}).

\SGAStatement{Remark 5.2.5}{The assertions of \SGARefLink{sga3:xvii:lemma:5:2:3:item:ii}{5.2.3(ii)} can also
be proved when \(U\) is not commutative, but we shall not need
this in order to prove \SGARefLink{sga3:xvii:theorem:5-1-1}{5.1.1}.}
\SGARefTarget{sga3:xvii:remark:5:2:5}

\SGASubsection{5.3}{\texorpdfstring{The Case Where \(H\) Is
Smooth}{The Case Where H Is Smooth}}
\SGARefTarget{sga3:xvii:subsection:5:3}

\SGAStatement{Proposition 5.3.1}{The assertions of \SGARefLink{sga3:xvii:lemma:5:2:3:item:i}{5.2.3(i)}
remain true when the hypothesis ``\(H\) is étale'' is replaced by
``\(H\) is smooth.''}
\SGARefTarget{sga3:xvii:proposition:5:3:1}

Proceeding as in the proof of \SGARefLink{sga3:xvii:lemma:5:2:3:item:i}{5.2.3(i)}, we reduce to the case in
which \(U\) is also commutative.

Let \(N'\) be the set of positive integers prime to \(p\), ordered
by divisibility.  For every \(n\in N'\), the group \({}_nH\) is
étale, and the family of groups \({}_nH\), \(n\in N'\), is
schematically dense in \(H\), since \(H\) is smooth
(\SGARefLink{sga3:ix:remark:4:10}{Exposé IX, 4.10}).  Let \(E_n\) be the inverse image of
\({}_nH\) in \(E\), so that \(E_n\) is an extension of
\({}_nH\) by \(U\), and let \(R_n\) be the \(k\)-functor of
lifts of \({}_nH\) in \(E_n\) (cf.\ \SGARefLink{sga3:xvii:lemma:5:2:3:item:ii}{5.2.3(ii)}).  If \(n\)
divides \(m\), there is plainly a natural \(k\)-morphism
\[
  R_m\longrightarrow R_n,
\]
so the \(R_n\) form an inverse system of \(k\)-functors.
